\def\stat{grusho+7}





\def\tit{ОБ АНАЛИЗЕ ОШИБОЧНЫХ СОСТОЯНИЙ В~РАСПРЕДЕЛЕННЫХ ВЫЧИСЛИТЕЛЬНЫХ СИСТЕМАХ$^*$}





\def\titkol{Об анализе ошибочных состояний 


в~распределенных вычислительных системах}








\def\aut{А.\,А.~Грушо$^1$, М.\,И.~Забежайло$^2$, А.\,А.~Зацаринный$^3$, 


А.\,В.~Николаев$^4$, В.\,О.~Писковский$^5$, В.\,В.~Сенчило$^6$, 


И.\,В.~Судариков$^7$, Е.\,Е.~Тимонина$^8$}





\def\autkol{А.\,А.~Грушо, М.\,И.~Забежайло, А.\,А.~Зацаринный и др.} 


%А.\,В.~Николаев,


%$^4$, В.\,О.~Писковский$^5$, В.\,В.~Сенчило$^6$,  И.\,В.~Судариков$^7$, Е.\,Е.~Тимонина$^8$}





\titel{\tit}{\aut}{\autkol}{\titkol}





\index{Грушо А.\,А.}


\index{Забежайло М.\,И.}


\index{Зацаринный А.\,А.}


\index{Николаев А.\,В.}


\index{Писковский В.\,О.}


\index{Сенчило В.\,В.}


\index{Судариков И.\,В.}


\index{Тимонина Е.\,Е.}


\index{Grusho A.\,A.}


\index{Zabezhailo M.\,I.}


\index{Zatsarinnyy A.\,A.}


\index{Nikolaev A.\,V.}


\index{Piskovski V.\,O.}


\index{Senchilo V.\,V.}


\index{Sudarikov I.\,V.}


\index{Timonina E.\,E.}





{\renewcommand{\thefootnote}{\fnsymbol{footnote}}


\footnotetext[1]


{Работа выполнена при частичной поддержке РФФИ (проект 15-29-07981~офи-м).}}





\renewcommand{\thefootnote}{\arabic{footnote}}


\footnotetext[1]{Институт проблем информатики Федерального исследовательского центра 


<<Информатика и~управ\-ле\-ние>> Российской академии наук, 


\mbox{grusho@yandex.ru}}


\footnotetext[2]{Институт проблем информатики Федерального исследовательского центра 


<<Информатика и~управ\-ле\-ние>> Российской академии наук,   


\mbox{m.zabezhailo@yandex.ru}}


\footnotetext[3]{Институт проблем информатики Федерального исследовательского центра 


<<Информатика и~управ\-ле\-ние>> Российской академии наук, 


\mbox{AZatsarinny@ipiran.ru}}


\footnotetext[4]{Институт химической физики им.\ 


Н.\,Н.~Семёнова Российской академии наук,\newline \mbox{gentoorion@mail.ru}}


\footnotetext[5]{Институт проблем информатики Федерального исследовательского центра 


<<Информатика и~управ\-ле\-ние>> Российской академии наук,


\mbox{vpvp80@yandex.ru}}


\footnotetext[6]{Институт проблем информатики Федерального исследовательского центра 


<<Информатика и~управ\-ле\-ние>> Российской академии наук,


\mbox{volodias@mail.ru}}


\footnotetext[7]{Институт проблем информатики Федерального исследовательского центра 


<<Информатика и~управ\-ле\-ние>> Российской академии наук,


\mbox{4seev3@gmail.com}}


\footnotetext[8]{Институт проблем информатики Федерального исследовательского центра 


<<Информатика и~управ\-ле\-ние>> Российской академии наук,


\mbox{eltimon@yandex.ru}}








 \label{st\stat}





                  \thispagestyle{headings}


                  


  \vspace*{-12pt}


  


  


 


  \Abst{Приведен анализ ошибочных состояний распределенных вычислительных систем на 


примере OpenStack, вызванных отсутствием синхронизации между компонентами облачной 


вычислительной среды (ОВС), а~также рассмотрены вероятные источники возникновения 


рассинхронизации. Пред\-став\-лен разбор журнальных записей при анализе одного из 


перечисленных ошибочных со\-сто\-яний программной платформы OpenStack. Показан метод 


анализа подобных ситуаций. Предложено построение и~поддержание в~актуальном со\-сто\-янии 


информационной модели управ\-ля\-емой ОВС. Обоснована 


не\-об\-хо\-ди\-мость использования такой модели для выполнения требований надежного управ\-ле\-ния 


ОВС, предоставления качественных облачных услуг типа сервис, 


платформа, инфраструктура, сетевые функции, цепочки сетевых функций.}


  


  \KW{распределенные вычислительные системы; облачные вы\-чис\-ли\-тель\-ные среды; 


OpenStack; RabbitMQ; Nova; Neutron; анализ ошибочных со\-сто\-яний; информационные облачные 


сетевые инфраструктуры}





\DOI{10.14357\!/\!08696527180108}








\vskip 10pt plus 9pt minus 6pt








      \thispagestyle{headings}


      


      \vspace*{-18pt} 





\section{Введение}





\vspace*{-6pt}





Некоторые свойства инфраструктуры распределенных 


ОВС порождают встречающиеся при их эксплуатации ошибочные 


состояния~[1]. 





  В~[1] обосновано, что одним из факторов, инициирующих сетевые сбои 


и~ошибки, является неизбежная рассинхронизация компонентов сети. 


Рас\-смот\-рим, как влияет физическое удаление инфраструктурных компонентов 


друг от друга на теоретическую возможность существования уязвимости 


в~инфраструктуре, вызванную естественной разницей системных часов, 


принадлежащих разным сетевым компонентам. Для простоты будем считать 


тактовую частоту обобщенного устройства равной~3~ГГц, что соответствует для 


оптоволокна характерному расстоянию $l\hm = 7$~см (показатель преломления 


принят равным~1,5). Для уверенной синхронизации взаимодействующие 


устройства должны находиться друг от друга на расстояниях, сравнимых с~$l$. 


В~противном случае при реконфигурации сети необходимо предусматривать 


дополнительные средства, обеспечивающие целостность и~непротиворечивость 


работы сети. Например, можно использовать методы, описанные в~работе~[1]. 





Возможно использование результатов работы~[2] и~деление сети на так 


называемые слои виртуальных подсетей. В~подсетях возможно независимое, 


полное и~непротиворечивое конфигурирование и~функционирование виртуальных 


машин (ВМ) и~сетевых компонентов, принадлежащих слою. Отметим, что пока удается 


найти только достаточные условия непротиворечивости работы сети. 


{\looseness=1





}


  


  Противоречия конфигураций приводят к~появлению уязвимостей в~течение 


временн$\acute{\mbox{о}}$го промежутка, который определяется временн$\acute{\mbox{ы}}$ми интервалами 


прохождения пакета между сетевыми компонентами. Можно ввести понятие 


поверхности атаки как отношение времени рассинхронизации двух серверов ко 


времени прохождения IP-па\-ке\-та между ними. Чем больше поверхность атаки, 


тем выше вероятность потенциального нарушителя внедрить свой код 


в~атакуемый узел сети. Отметим, что возможности такой атаки возрастают, если 


атакующее программное обеспечение (ПО) находится физически ближе к~цели, чем контролирующий 


конфигурацию сети сервер (по аналогии с~компьютерными играми и~задержкой 


реакции игры на воздействие игрока). 


  


  Для контролируемой защищенной работы инфраструктуры необходимо, чтобы 


корректирующее воздействие на узлы и~компоненты проходило быстрее 


воздействия атакующей стороны на те же компоненты. А~это значит, что 


управление распределенной ОВС, включая принятые меры противодействия 


вирусам, должно физически находиться как можно ближе к~критическим 


компонентам инфраструктуры. 





\section{Ошибки синхронизации в~распределенных облачных вычислитедьных


средах на~примере 


OpenStack}





  Рассмотрим помимо описанного ранее <<нештатного>> поведения компонент 


OpenStack из-за нехватки ресурсов RabbitMQ~[3] другие подобные ошибки 


в~комплексе ПО OpenStack, которые представлены в~таблице.


  


  \begin{table}\small


  \begin{center}


  


  \begin{tabular}{|p{30mm}|p{91mm}|}


  \multicolumn{2}{c}{Ошибки в~комплексе ПО OpenStack}\\


  \multicolumn{2}{c}{\ }\\[-6pt]


  \hline


\multicolumn{1}{|c|}{Ошибка}&\multicolumn{1}{c|}{Описание}\\


\hline


1.~Нехватка ресурсов RabbitMQ приводит к~потере сетевого соединения, распад облачной 


структуры на <<подоблака>>~\cite{3-g7}&Нехватка числа файловых дескрипторов 


и~последующей перезагрузки системы доставки сообщений RabbitMQ\\


\hline


2.~Обрыв соединения со службой RabbitMQ во время миграции ВМ приводит 


к~сбою в~машине состояний Nova~\cite{4-g7}&Обрыв соединения с~упомянутой службой 


RabbitMQ во время миграции ВМ приводит к~сбою в~машине состояний 


Nova. При этом статусом в~БД Nova остается <<миграция>> и~указание на узел, 


с~которого мигрировали ВМ, хотя ВМ работают на целевом узле. Эта ошибка, скорее всего, имеет 


своими корнями некорректную обработку сложного дерева состояний Nova, возникающего при 


обрыве связи со службой доставки сообщений. Скорее всего, вместо корректной обработки всех 


возможных состояний в~коде Nova используется просто набор тайм-аутов, не обновляя и~не 


синхронизируя актуальные состояния\\


\hline


3.~При перезапуске ВМ сетевое подключение не формируется заново и~ВМ остается без 


доступа к~сети~\cite{5-g7}&Если первая попытка запуска ВМ по каким-то причинам не 


увенчалась успехом, но при этом сетевое подключение на запущенной ВМ не было уничтожено, 


то при перезапуске ВМ сетевое подключение не формируется заново и~ВМ остается без доступа к~сети. Данная ошибка, скорее всего, обусловлена некорректной обработкой сбоя при создании 


ВМ в~сетевой компоненте Neutron. При этом если у ВМ имеются два сетевых интерфейса, то из 


альтернативной сети успешный запуск возможен \\


\hline


4.~Команда безусловного уничтожения ВМ, отданная на этапе создания блочного устройства, 


приводит к~сбою~\cite{6-g7}&В распределенных сис\-те\-мах, компоненты которых 


разрабатываются без жестко зафиксированного единого проекта, фактически невозможно 


полностью корректно обработать (и~даже выписать) все возможные промежуточные 


со\-сто\-яния, в~которых может оказаться рассматриваемая промежуточная сис\-те\-ма. В данной ошибке команда 


безусловного уничтожения ВМ, отданная на этапе создания блочного устройства, приводит 


к~сбою\\


\hline


5.~Ошибка при удалении текущего IP-ад\-ре\-са~\cite{7-g7}&Ошибка может возникнуть при 


нарушении целостности актуальной информации об объектах тенанта при конкурентном 


выполнении команд <<удалить ВМ>> и~<<удалить текущий IP-ад\-рес>> \\


\hline


6.~Ошибка миграции при наличии у~ВМ дополнительного диска~\cite{8-g7}&Происходит при 


<<живой>> миграции ВМ под управлением Nova (libvirt), име\-ющей присоединенное блочное 


устройство на cinder на основе Ceph. Вероятнее всего, данная ошибка возникает из-за 


несогласованности API Nova с~поддержкой libvirt операций с~Ceph remote block device (rbd)\\


\hline


\multicolumn{2}{r}{\textit{Окончание таблицы на с.}~\pageref{e1}}


\end{tabular}


\end{center}


\end{table}





  


  


 \pagebreak


 


  Указанные выше уязвимости имеют несколько общих черт: 


  \begin{itemize}


\item не завершена разработка и~не соблюдается протокол взаимодействия 


компонентов;


\item не стабилизирована работа программных интерфейсов и~не 


соблюдаются требования к~их использованию.


\end{itemize}





\begin{table}\small


  \begin{center}


  \label{e1}


  


  \begin{tabular}{|p{30mm}|p{91mm}|}


  \multicolumn{2}{c}{Ошибки в~комплексе ПО OpenStack (\textit{окончание})}\\


  \multicolumn{2}{c}{\ }\\[-6pt]


  \hline


\multicolumn{1}{|c|}{Ошибка}&\multicolumn{1}{c|}{Описание}\\


\hline


7.~Удаленные пользователи OpenStack могут обходить защиту от подмены IP, вклиниваясь 


в~процедуру запуска ВМ тенанта~\cite{9-g7}&Уязвимость CVE-2015-5240 (Race 


Conditions vulnerability in Openstack Neutron) обусловлена тем, что число состояний алгоритмов, 


использующих распределенные компоненты, очень велико и~корректная обработка 


разнообразных нештатных ситуаций в~промежуточных состояниях является скорее 


исключением, чем правилом\\


\hline


8.~Обращение в~БД к~записи, которая еще не завершена~\cite{10-g7} (Аномалия~2. P1.Dirty 


Read)


&Имеет место при первом запуске ВМ на сервисе Nova, когда в~БД планировщика 


еще нет полной информации о~дисковых ресурсах этого сервиса. По-видимому, это связано 


с~рассогласованностью формата БД планировщика и~инсталляционных 


и~инициализационных процедур сервиса Nova: планировщик пытается прочитать в~БД запись, 


которая еще полностью не завершена \\


\hline


9.~Задержки при работе с~БД~\cite{11-g7}&В некоторых таблицах, используемых 


компонентом Nova, продублированы индексы, что приводит к~неоправданным задержкам при 


работе с~такими таблицами \\


\hline


10.~Попытка присоединить дополнительный том для хранения данных к~экземпляру ВМ 


игнорируется~\cite{12-g7}&Невозможно присоединить и~использовать дополнительный заново 


созданный том для хранения данных, причем этот том отображается как доступный для работы. 


Причина, по-ви\-ди\-мо\-му, в~неполноте информации, доступной для работы Nova, которая,


 в~свою 


очередь, не может выполнить команду QEMU по причине дублирования\\


\hline


11.~Невозможно удалить образ ВМ, с~которой запущен экземпляр ВМ~\cite{13-g7}&В БД 


хранится свойство образа, которое указывает на факт работающей ВМ, созданной из образа. 


При попытке удаления образа информация об этом указателе не сбрасывается и~не дает удалить 


образ вопреки заявленной функциональности\\


\hline


  \end{tabular}


  \end{center}


  \end{table}





  При работе компонентов не соблюдаются принятые нормы и~правила при 


обработке исключений и~ошибок. Возможно, это недостатки проектирования. 


  


  Традиционно обнаружение подобных ошибок осуществляется во время 


эксплуатации. Как правило, при обнаружении какой-либо аномалии 


в~функционировании системы осуществляется детальное изучение журналов 


событий и~сопоставление ошибок и~предупреждений. Фактически первичный 


анализ сводится к~установлению корреляций между сообщениями в~журналах 


различных подсистем. Затем осуществляется поиск подобных ошибок 


в~Интернете и~при нахождении таковых проводится анализ причин их 


возникновения. При возможности, осуществляется <<увязка>> по схеме  


<<при\-чи\-на--след\-ст\-вие>> пар сообщений об ошибках. Далее формируется 


первичная <<модель>> последовательности событий, которая проверяется на 


тестовых примерах. Тестовые примеры затем используются как <<шаги 


воспроизведения>> в~описании уже новой ошибки в~базе данных (БД) системы 


контроля ошибок и~управления циклом жизни~ПО. 


{ %\looseness=1





}


  


  В то же время, например, ошибка~3 (см.\ таблицу)~\cite{5-g7}, связанная с~неправильным 


сетевым подключением перезапущенной ВМ в~относительно современном 


окружении OpenStack release Mitaka, не может быть проанализирована 


стандартными средствами. Суть ошибки заключается в~том, что ВМ пересоздается 


на Nova, отличной от первоначальной, а~сетевой порт остается на агенте 


компонента Neutron на том же хосте, что и~первоначальная Nova. При этом 


стандартные средства Neutron позволяют обнаружить <<привязку>> порта 


к~физическому хосту только по идентификатору ВМ и~хосту ее расположения, 


который после перезапуска ВМ уже изменился. Соответственно, чтобы просто 


увязать факты, что агент поддерживает старый сбойной порт с~присутствием на 


хосте, на котором он реализован, надо анализировать либо журналы Nova 


и~Neutron, собранные со всех Nova-хостов облака, либо конфигурационные базы 


OpenStack вручную. Процесс это не быстрый, а~в~интенсивно загруженной среде 


и~без централизованного упорядоченного журнала событий, возможно, еще 


и~дающий неоднозначный результат. 


  


  Далее, обнаружив хост локализации сбойного порта, придется увязать записи из 


журнала агента Neutron с~журналом утилиты, которая проводила развертывание 


ВМ и~виртуальной сети, например Heat. Убедившись в~соответствии временн$\acute{\mbox{ы}}$х 


меток и~идентификаторов запросов, порта, ВМ и~анализируемой подсети, можно 


будет приступить к~попытке моделировать подобную ситуацию. Например, можно 


выделить отдельную зону с~неправильно сконфигурированной Nova, которая 


будет неспособна развернуть ВМ и~пытаться воспроизвести ситуацию 


с~перезапуском ВМ на другой Nova. 


  


  Очевидно, что стандартными средствами такую задачу также решить 


невозможно и~придется прибегать к~некоторым ухищрениям, тщательно следя за 


тем, чтобы сами эти приемы не внесли своего вклада в~особенности генерируемой 


ошибки. В~описании данной ошибки для ее воспроизведения предлагается 


модифицировать код Nova на тестовом узле, что сложно признать хорошим 


методом тестирования по любым канонам. 


   


   Разберем детально журнальные записи для ситуации, при которой возникает 


ошибка~5 при удалении текущего IP-ад\-ре\-са (см.\ таблицу)~\cite{7-g7}. В~описании указано, что 


при манипуляции с~ВМ и~созданием, остановкой и~удалением 


возникает необходимость удаления ассоциированного с~ВМ IP-ад\-ре\-са, которое 


иногда приводит к~возникновению ошибки, связанной с~тем, что собственно 


<<текущий IP-ад\-рес>> к~моменту исполнения команды может уже не 


существовать в~результате выполнения конкурентно запущенной команды по 


удалению ВМ. Ошибка относится к~так называемому <<мерцающему>> типу, 


когда факт ее возникновения зависит от массы характеристик слабо 


контролируемого окружения. Такой тип ошибок является, пожалуй, одним из 


самых сложных и~неприятных из всех возможных для разработчика. Чтобы 


определить, что же послужило причиной ошибки, час\-то приходится проводить 


массу экспериментов, обрабатывать огромный объем генерируемых журнальных 


записей и~данных.


   


   Ручной разбор журнальных записей достаточно трудоемкий в~силу того, что 


приходится кроме упорядочивания событий по времени восстанавливать скрытые 


взаимосвязи и~ассоциации между объектами. Отсутствие актуальной модели ОВС, 


объединяющей описание как отдельных объектов инфраструктуры, так и~связей 


между ними, неизбежно приводит не только к~нарушению це\-лост\-ности 


конфигурации, но и~к~появлению ошибок управления ОВС, в~том чис\-ле 


<<мер\-ца\-ющих>>. Методологическая причина этого кроется в~нарушении порядка 


и~синхронизации процессов управления~\cite{1-g7, 3-g7}. 





Выход из ситуации 


видится в~решении задачи создания и~поддержания в~актуальном состоянии 


информационной модели ап\-па\-рат\-но-про\-грам\-мно\-го комплекса, включая 


конфигурацию виртуальной среды ОВС. В~традиционном варианте реализации  


ап\-па\-рат\-но-про\-грам\-мных комплексов эта задача успешно решается 


средствами класса IBM Rational\,/\,Telelogic Tau~\cite{14-g7, 15-g7}. 


В~традиционном подходе применение инструментов подобного класса (IBM 


Rational\,/\,Telelogic DOORS/Tau/Synergy/TauTester/DocExpress) позволяет 


значительно оптимизировать выполнения ряда ITIL


(Information Technology Infrastructure Library) про\-цес\-сов, в~частности 


Change и~Release Management. В~случае же ОВС, когда характерные времена 


уменьшаются до секунд, наличие подобного отлаженного инструментария 


является одним из необходимых условий проведения успешной реконфигурации 


и~защиты от нежелательных вторжений. 


   


   Применение подобных средств даст возможность:


   \begin{itemize}


\item упорядочить процессы управления ОВС, сохраняя целостность 


конфигурации ОВС в~практически любой момент времени; 


\item определять временн$\acute{\mbox{ы}}$е периоды и~управлять ими, когда ОВС не 


находится в~состоянии, отвечающем требованиям SLA (Service


Level Agreement) и~политик 


безопасности, и~поэтому уязвима и~ограниченно работо\-спо\-собна;


{\looseness=1





}





\item декомпозировать ОВС на составляющие компоненты, тем самым 


понижая размерность и~сложность как каждого компонента, так и~ОВС 


в~целом, а значит, переводя задачу верификации процессов реконфигурации 


ОВС в~разряд разрешимых. 


\end{itemize}





\section{Выводы}


 


   Подводя итог приведенных выше рассуждений, можно заключить, что для 


стабилизации работы продукта необходимо: 





\pagebreak





\noindent


   \begin{itemize}


\item предусмотреть центр регистрации ошибок и~мониторинга;


\item конфигурацию ОВС проводить в~условиях обеспечения 


мультивариантности, процесс реконфигурации должен быть 


двухфазным~\cite{1-g7}; 


\item план реконфигурации должен быть по возможности правильным 


и~непротиворечивым.


\end{itemize}





   С этой целью граф модели ОВС должен состоять из подмножеств 


плоскостей~\cite{2-g7}, оптимально обеспечивающих правильность, 


непротиворечивость и~упорядоченность плана реконфигурации~\cite{1-g7}. 


Конфигурации должны регистрироваться в~БД с~поддержкой транзакций и~двухфазного коммита. Уровень изоляции транзакций зависит от необходимого 


уровня обслуживания и~кри\-тич\-ности работы компонентов. Регистрация событий 


долж\-на вестись в~централизованном иерархическом порядке в~соответствии 


с~планом реконфигурации ресурса. Для регистрации и~ситуационного анализа 


работы ОВС, а~также для возможности определить время, узел, характер и~вектор 


атаки вредоносного ПО, если таковая будет иметь место. Помимо смысловых 


меняющихся полей необходимо регистрировать: 


   \begin{itemize}


\item отметки времени узлов, занятых в~транзакциях, включая БД и~серверы 


регистрации событий (журналирование). В~случае иерархии таких серверов, 


при анализе данных необходимы отметки времени серверов, участвующих 


в~регистрации; 


\item идентификатор версии конфигурации; 


\item если ОВС выполнена с~использованием SDN (Software-Defined Networking), 


идентификатор 


<<плоскости>>, slice, в~которой проводилась реконфигурация; 


\item идентификаторы проекта, тенанта, объекта конфигурации и~их 


взаимосвязи.


\end{itemize}





Под объектом конфигурации здесь понимается сущность, влияющая на 


работу инфраструктуры, как-то: узел (физический и~виртуальный), IP-ад\-рес, 


порт, протокол, приложение, сервис, сетевая функция, сформированные 


цепочки сетевых функций, учетная запись, права доступа и~т.\,д.


   


   К сожалению, на настоящий момент не существует ни теоретически 


обоснованных подходов к~построению целостных непротиворечивых 


конфигураций в~распределенных ОВС, ни целостных методов автоматической 


проверки, способных адекватно тестировать все состояния распределенных 


систем. Соответственно, ошибки в~них выявляются по большей части на этапе 


эксплуатации. Естественно ожидать, что даже после длительного цикла 


эксплуатации и~исправления ошибок все равно в~системе будут оставаться 


проблемы, способные привести к~ошибочным состояниям при маловероятных 


стечениях обстоятельств. Наличие поддерживаемой в~актуальном состоянии 


информационной модели ап\-па\-рат\-но-про\-грам\-мно\-го комплекса ОВС позволит 


существенно сократить путь к~решению задачи управления и~поддержки 


целостных непротиворечивых конфигураций в~распределенных ОВС.


   


{\small\frenchspacing


{%\baselineskip=10.8pt


\addcontentsline{toc}{section}{Литература}


\begin{thebibliography}{99}


\bibitem{1-g7}


\Au{Грушо А.\,А., Забежайло М.\,И., Зацаринный~А.\,А., Писковский~В.\,О.} Безопасная 


автоматическая реконфигурация облачных вычислительных сред~/\!\!/ Системы и~средства 


информатики, 2016. Т.~26. №\,3. С.~83--92.


\bibitem{2-g7}


\Au{McGeer R.} Declarative verifiable SDI specifications~/\!\!/ 37th IEEE Symposium on Security 


and Privacy Proceedings.~--- IEEE, 2016. P.~198--203. {\sf http:// spw16.langsec.org/papers/mcgeer-


verifiable-sdi-specs.pdf}.





\bibitem{3-g7}


\Aue{Грушо А.\,А., Забежайло М.\,И., Зацаринный~А.\,А., Николаев~А.\,В., Писковский~В.\,О., 


Тимонина~Е.\,Е.} Классификация ошибочных состояний в~распределенных вычислительных 


системах и~источники их возникновения~/\!\!/ Системы и~средства информатики, 2017. Т.~27. 


№\,2. С.~29--40. 


\bibitem{4-g7}


OpenStack Compute (nova): Inconsistent state when connection to conductor is lost during live 


migration, 2016. {\sf https://bugs.launchpad.net/nova/+bug/1536589}.


\bibitem{5-g7}


OpenStack Compute (nova): Network not always cleaned up when spawning VMs, 2016. {\sf 


https://bugs.launchpad.net/nova/+bug/1597596}. 


\bibitem{6-g7}


OpenStack Compute (nova): Nova force-delete can't delete instance in vm\_state 


block\_device\_mapping, 2017. {\sf https://bugs.launchpad.net/nova/+bug/1704945}. 


\bibitem{7-g7}


OpenStack Compute (nova): Race conditions between compute and schedule disk report. {\sf 


https://bugs.launchpad.net/nova/+bug/1704975}.


\bibitem{8-g7}


OpenStack Compute (nova): Live migration fails with an attached non-bootable Cinder volume (Pike),


2017. 


{\sf https://bugs.launchpad.net/nova/+bug/1715569}. 


\bibitem{9-g7}


Neutron: Race Conditions vulnerability in Openstack Neutron, 2015. {\sf  


https:// cyber.vumetric.com/vulns/CVE-2015-5240/race-conditions-vulnerability-openstack-neutron}.


\bibitem{10-g7}


OpenStack Compute (nova): Race conditions between compute and schedule disk report,


2016. {\sf 


https://bugs.launchpad.net/nova/+bug/1610679}.


\bibitem{11-g7}


OpenStack Compute (nova): Duplicate indexes in nova-db, 2016. {\sf 


https:// bugs.launchpad.net/nova/+bug/1641185}.


\bibitem{12-g7}


OpenStack Compute (nova): Can't attach volume to volume-backed instance, 2017. {\sf 


https://bugs.launchpad.net/nova/+bug/1678694}, 2017.


\bibitem{13-g7}


OpenStack Compute (nova): Can not delete the image has been launched instance when use rbd. {\sf 


https://bugs.launchpad.net/nova/+bug/1697391}.


\bibitem{14-g7}


\Au{De Wet N., Kritzinger~P.} Towards model-based communication protocol performability analysis 


with UML~2.0~/\!\!/ Southern African Telecommunication Networks and Applications Conference 


(SATNAC) Proceedings.~--- Spier Wine Estate, 2004. 6~р. {\sf 


http://www.satnac.org.za/proceedings/2004/Software/No\%20105\%20-\%20de\%20Wet.pdf}.


\bibitem{15-g7}


Standards-based, model-driven development solution for complex systems, 2016. {\sf  


https://www-03.ibm.com/software/products/en/ratitau}.





        \end{thebibliography}


} }








\vspace*{-3pt}





\hfill{\small\textit{Поступила в~редакцию 04.02.18}}














%\vspace*{8pt}





%\hrule





%\vspace*{2pt}





\newpage





\vspace*{-28pt}





%\hrule





%\vspace*{9pt}











\def\tit{ABOUT THE~ANALYSIS OF~ERRATIC STATUSES 


IN~THE~DISTRIBUTED COMPUTING SYSTEMS}





\def\titkol{About the~analysis of~erratic statuses 


in~the~distributed computing systems}








\def\aut{A.\,A.~Grusho$^1$, M.\,I.~Zabezhailo$^1$, A.\,A.~Zatsarinny$^1$, 


A.\,V.~Nikolaev$^2$, V.\,O.~Piskovski$^1$, V.\,V.~Senchilo$^1$, 


I.\,V.~Sudarikov$^1$, and~E.\,E.~Timonina$^1$}


\def\autkol{A.\,A.~Grusho, M.\,I.~Zabezhailo, A.\,A.~Zatsarinnyy, et al.}


%A.\,V.~Nikolaev, 


%V.\,O.~Piskovski$^1$, V.\,V.~Senchilo$^1$,  I.\,V.~Sudarikov$^1$, and~E.\,E.~Timonina$^1$}











\titel{\tit}{\aut}{\autkol}{\titkol}





\vspace*{-9pt}








 \noindent


   $^1$Institute of Informatics Problems, Federal Research Center ``Computer 


Sciences\linebreak


$\hphantom{^1}$and Control'' of the Russian Academy of Sciences; 44-2~Vavilov Str., Moscow\linebreak 


$\hphantom{^1}$119133, Russian Federation


   


   \noindent


   $^2$N.\,N.~Semenov Institute of Chemical Physics, Russian Academy of Sciences;\linebreak


   $\hphantom{^1}$4-1~Kosygina Str.,  Moscow 119991, Russian Federation








\def\leftfootline{\small{\textbf{\thepage}


\hfill Sistemy i Sredstva Informatiki~---


Systems and Means of Informatics 2018 vol~28 no\,1}


}%


 \def\rightfootline{\small{Sistemy i Sredstva Informatiki~---


 Systems and Means of Informatics   2018   vol~28   no\,1


\hfill \textbf{\thepage}}}


    


    


    


   


   


   \Abste{The analysis of erratic statuses of the distributed computing systems on the example of 


OpenStack, caused by the absence of synchronization


between components of a cloudy computing 


environment and, also, probable sources of their origin are considered. 


Detail analysis of 


journal records is provided in the analysis of one of the listed erratic statuses of software platform of 


OpenStack. The method of the analysis of similar situations is shown. Creation and maintenance in 


current state of an information model of a managed cloudy computing environment is 


suggested. The necessity to use 


such model for execution of requirements of positive control of a cloudy computing 


environment and provision of high-quality cloudy services 


of such types as service, platform, 


infrastructure, network functions, chains of network functions is substantiated.} 





   \KWE{distributed computing systems; cloudy computing environments; OpenStack; 


RabbitMQ; Nova; Neutron; analysis of erratic statuses; information cloudy network infrastructures}


   





   


\DOI{10.14357\!/\!08696527180108}





%\vspace*{-24pt}





   \Ack


   \noindent


   The paper was partially supported by the


   Russian Foundation for Basic Research (project 15-29-07981~ofi-m).








\renewcommand{\bibname}{\protect\rmfamily References}


%\renewcommand{\bibname}{\large\protect\rm References}





{\small\frenchspacing


{%\baselineskip=10.8pt


\addcontentsline{toc}{section}{References}


\begin{thebibliography}{99}





\bibitem{1-g7-1}


\Aue{Grusho, A.\,A., M.\,I.~Zabezhailo, A.\,A.~Zatsarinnyy, and V.\,O.~Piskovski}. 2016. 


Bez\-opas\-naya avtomaticheskaya rekonfiguratsiya oblachnykh vychislitel'nykh sred [Secure automatic 


reconfiguration of cloudy computing]. \textit{Sistemy i~Sredstva Informatiki~--- Systems and Means of 


Informatics} 26(3):83--92.


\bibitem{2-g7-1}


\Aue{McGeer, R.} 2016. Declarative verifiable SDI specifications. \textit{37th IEEE 


Symposium on 


Security and Privacy Proceedings}. IEEE. 198--203. Available at: {\sf 


http:// spw16.langsec.org/papers/mcgeer-verifiable-sdi-specs.pdf} (accessed February~3, 2018).


\bibitem{3-g7-1}


\Aue{Grusho, A.\,A., M.\,I.~Zabezhailo, A.\,A.~Zatsarinnyy, A.\,V.~Nikolaev,  


V.\,O.~Piskovski, and E.\,E.~Timonina.} 2017. Klassifikatsiya oshibochnykh sostoyaniy 


v~raspredelennykh vychislitel'nykh sistemakh i~istochniki ikh vozniknoveniya [Erroneous states 


classification in distributed computing systems and sources of their occurrences]. \textit{Sistemy 


i~Sredstva Informatiki~--- Systems and Means of Informatics} 27(2):29--40. 


\bibitem{4-g7-1}


OpenStack Compute (nova): Inconsistent state when connection to conductor is lost during live 


migration.  2016. Available at: {\sf https://bugs.launchpad.net/nova/ +bug/1536589} (accessed Febrary~3, 


2018).


\bibitem{5-g7-1}


OpenStack Compute (nova): Network not always cleaned up when spawning VMs. 2016. Available at: 


{\sf https://bugs.launchpad.net/nova/+bug/1597596} (accessed February~3, 2018). 


\bibitem{6-g7-1}


OpenStack Compute (nova): Nova force-delete can't delete instance in vm\_state 


block\_device\_mapping. 2017. Available at: {\sf https://bugs.launchpad.net/nova/+bug/ 1704945} 


(accessed February~3, 2018). 


\bibitem{7-g7-1}


OpenStack Compute (nova): Race conditions between compute and schedule disk report. Available at: 


{\sf https://bugs.launchpad.net/nova/+bug/1704975. 2017}  (accessed February~3, 2018).


\bibitem{8-g7-1}


OpenStack Compute (nova): Live migration fails with an attached non-bootable Cinder volume (Pike). 


2017. Available at: {\sf https://bugs.launchpad.net/nova/+bug/1715569} (accessed February~3, 2018). 


\bibitem{9-g7-1}


Neutron: Race Conditions vulnerability in Openstack Neutron. 2015. Available at: {\sf 


https://cyber.vumetric.com/vulns/CVE-2015-5240/race-conditions-vulnerability-openstack-neutron/} (accessed 


February~3, 2018).


\bibitem{10-g7-1}


OpenStack Compute (nova): Race conditions between compute and schedule disk report. 2016. 


Available at: {\sf https://bugs.launchpad.net/nova/+bug/1610679} (accessed February~3, 2018).


\bibitem{11-g7-1}


OpenStack Compute (nova): Duplicate indexes in nova-db. 2016. Available at: {\sf 


https:// bugs.launchpad.net/nova/+bug/1641185} (accessed February~3, 2018).


\bibitem{12-g7-1}


OpenStack Compute (nova): Can't attach volume to volume-backed instance. 2017. Available at: {\sf 


https://bugs.launchpad.net/nova/+bug/1678694} (accessed February~3, 2018).


\bibitem{13-g7-1}


OpenStack Compute (nova): Can not delete the image has been launched instance when use rbd. 2017. 


Available at: {\sf https://bugs.launchpad.net/nova/+bug/1697391} (accessed February~3, 2018).


\bibitem{14-g7-1}


\Aue{De Wet, N. and P.~Kritzinger.} 2004. Towards model-based communication protocol 


performability analysis with UML~2.0. \textit{Southern African Telecommunication Networks and 


Applications Conference (SATNAC) Proceedings}. Spier Wine Estate. 6~р. Available at: {\sf 


http://www.satnac.org.za/proceedings/2004/Software/No\%20105\%20-\%20de\%20Wet.pdf} (accessed 


Febrary~3, 2018).


\bibitem{15-g7-1}


Standards-based, model-driven development solution for complex systems. 2015. Available at: {\sf 


https://www-03.ibm.com/software/products/en/ratitau} (accessed February~3, 2018).





\end{thebibliography}


} }








\vspace*{-3pt}





\hfill{\small\textit{Received February 4, 2018}}   





\Contr





\noindent


\textbf{Grusho Alexander A.} (b.\ 1946)~--- Doctor of Science in physics and 


mathematics, professor, Head of Laboratory, Institute of Informatics Problems, 


Federal Research Center ``Computer Sciences and Control'' of the Russian Academy 


of Sciences; 44-2~Vavilov Str., Moscow 119133, Russian Federation;  


\mbox{grusho@yandex.ru} 





\vspace*{3pt}





\noindent


\textbf{Zabezhailo Michael I.} (b.\ 1956)~--- Doctor of Science in physics and 


mathematics, Head of Laboratory, Institute of Informatics Problems, Federal 


Research Center ``Computer Sciences and Control'' of the Russian Academy of 


Sciences; 44-2~Vavilov Str., Moscow 119133, Russian Federation; 


\mbox{m.zabezhailo@yandex.ru} 





\vspace*{3pt}





\noindent


\textbf{Zatsarinny Alexander A.} (b.\ 1951)~--- Doctor of Science in technology, 


professor, Deputy Director, Federal Research Center ``Computer Sciences and 


Control'' of the Russian Academy of Sciences; 44-2~Vavilov Str., Moscow 119133, 


Russian Federation; \mbox{AZatsarinny@ipiran.ru} 





\vspace*{3pt}





\noindent


\textbf{Nikolaev Andrei V.} (b.\ 1973)~--- Candidate of Science (PhD) in physics 


and mathematics, senior scientist, N.\,N.~Semenov Institute of Chemical Physics,  


Russian Academy of Sciences; 4-1~Kosygina Str., Moscow 119991, Russian 


Federation; \mbox{gentoorion@mail.ru} 





\vspace*{3pt}





\noindent


\textbf{Piskovski Viktor O.} (b.\ 1963)~--- Candidate of Science (PhD) in physics 


and mathematics, senior scientist, Institute of Informatics Problems, Federal 


Research Center ``Computer Sciences and Control'' of the Russian Academy of 


Sciences; 44-2~Vavilov Str., Moscow 119133, Russian Federation; 


\mbox{vpvp80@yandex.ru}





\vspace*{3pt}





\noindent


\textbf{Senchilo Vladimir V.} (b.\ 1963)~--- scientist, Institute of Informatics 


Problems, Federal Research Center ``Computer Sciences and Control'' of the Russian 


Academy of Sciences; 44-2~Vavilov Str., Moscow 119133, Russian Federation; 


\mbox{volodias@mail.ru} 





\vspace*{3pt}





\noindent


\textbf{Sudarikov Igor V.} (b.\ 1989)~--- scientist, Institute of Informatics 


Problems, Federal Research Center ``Computer Sciences and Control'' of the Russian 


Academy of Sciences; 44-2~Vavilov Str., Moscow 119133, Russian Federation; 


\mbox{4seev3@gmail.com} 





\vspace*{3pt}





\noindent


\textbf{Timonina Elena E.} (b.\ 1952)~---


Doctor of Science in technology, professor, leading scientist, Institute of Informatics 


Problems, Federal Research Center ``Computer Sciences and Control'' of the Russian 


Academy of Sciences; 44-2~Vavilov Str., Moscow 119133, Russian Federation; 


\mbox{eltimon@yandex.ru} 





\label{end\stat}








\renewcommand{\bibname}{\protect\rm Литература}


